\section{Chapitre 4 : Intelligence Artificielle et Mécanismes de Proctoring}
\addcontentsline{toc}{section}{Chapitre 4 : Intelligence Artificielle et Mécanismes de Proctoring}

L'intelligence artificielle représente l'épine dorsale technologique qui confère à la plateforme Fraudly son qualificatif de système d'apprentissage agentique, la distinguant radicalement des dépôts de contenus numériques statiques d'antan. Au cœur de cette prouesse cognitive se trouvent les grands modèles de langage, dits Large Language Models, qui constituent de vastes réseaux de neurones artificiels dont l'architecture est basée sur le mécanisme d'attention introduit par le paradigme Transformer. La sélection du modèle cognitif optimal a fait l'objet d'une phase de recherche et de développement rigoureuse, matérialisée par un rapport comparatif exhaustif évaluant la vélocité d'inférence, la précision sémantique et la viabilité économique de multiples candidats. Les investigations ont conduit à retenir le modèle LLaMA 3 dans sa version paramétrique intermédiaire, plébiscité pour son excellence dans la génération de la langue française, sa capacité à produire des raisonnements structurés et sa latence d'exécution, qui s'est avérée drastiquement inférieure au seuil critique de deux secondes exigé pour les interactions conversationnelles fluides en temps réel. Pour s'affranchir d'une dépendance monolithique envers un fournisseur technologique unique et pour orchestrer harmonieusement les interactions complexes avec le modèle de fondation, l'architecture de Fraudly intègre la bibliothèque d'abstraction LangChain. Cette couche d'intermédiation logicielle permet de concevoir des chaînes de traitement conceptuelles sophistiquées, reliant le modèle de langage à divers outils périphériques, à l'historique conversationnel et à la mémoire vectorielle, tout en offrant la flexibilité de basculer instantanément vers un modèle alternatif de secours en cas de défaillance de l'interface de programmation principale.

L'une des innovations majeures de la plateforme réside dans l'implémentation du paradigme de Génération Augmentée par la Recherche, universellement connu sous l'acronyme RAG, qui propulse l'agent tuteur pédagogique. L'inconvénient inhérent aux modèles de langage réside dans leur propension à générer des informations factuellement erronées, qualifiées d'hallucinations, lorsqu'ils sont confrontés à des questions pointues dépassant les limites de leurs données d'entraînement initiales. La technique d'augmentation par la recherche pallie brillamment cette lacune en conditionnant strictement la phase de génération de texte à un sous-ensemble de connaissances externes préalablement vérifiées. Lorsqu'un étudiant adresse une requête conceptuelle au tuteur virtuel, le système ne se contente pas d'interroger directement le modèle de langage. Dans une première phase, la requête de l'étudiant est mathématiquement convertie en un vecteur d'enchâssement multidimensionnel représentant sa signature sémantique. Ce vecteur est ensuite confronté à l'ensemble des vecteurs indexés dans la base de données vectorielle ChromaDB, qui contient l'exhaustivité des contenus pédagogiques préalablement segmentés du cours concerné. Les algorithmes de recherche de proximité spatiale identifient alors les paragraphes du cours dont la signification est la plus connexe à la question posée. Dans la seconde phase du processus, le modèle de langage est invoqué à l'aide d'une instruction contraignante lui intimant de formuler une réponse pédagogique, claire et concise, en s'appuyant exclusivement et rigoureusement sur le contexte textuel extrait lors de la phase de recherche. Ce mécanisme garantit que les explications fournies à l'apprenant sont non seulement justes d'un point de vue académique, mais qu'elles s'inscrivent parfaitement dans le lexique, la progression et la méthodologie spécifiques adoptés par l'enseignant de la matière concernée.

Parallèlement à l'assistance pédagogique, le module d'intelligence artificielle déploie toute sa puissance analytique au service de la préservation de l'intégrité des évaluations à distance, un défi majeur abordé par le sous-système de télésurveillance, ou proctoring. Contrairement aux approches de surveillance traditionnelles reposant sur la présence synchrone de surveillants humains tentant de scruter simultanément des dizaines de flux vidéo, Fraudly automatise entièrement la détection des anomalies comportementales et environnementales grâce à la vision par ordinateur et à l'analyse algorithmique des interactions numériques. Ce mécanisme de télésurveillance est intrinsèquement multimodal. Le premier vecteur d'analyse repose sur le traitement en temps réel des trames vidéo capturées par la caméra de l'étudiant. Ces images sont soumises en continu à un algorithme de détection d'objets de l'état de l'art, appartenant à la famille des architectures You Only Look Once, communément désignée sous l'acronyme YOLO. Entraîné sur des milliers d'images, ce réseau de neurones convolutif possède l'acuité visuelle nécessaire pour repérer quasi instantanément la présence furtive de dispositifs prohibés dans le champ de vision, tels que des téléphones cellulaires, des tablettes électroniques ou des écouteurs dissimulés. Simultanément, des algorithmes spécialisés dans la reconnaissance faciale, tels que les réseaux convolutifs en cascade, s'assurent non seulement que le visage détecté correspond bien au gabarit biométrique de l'étudiant légitimement inscrit, mais ils vérifient également que l'apprenant est l'unique individu présent devant l'écran, levant une alerte immédiate si un tiers pénètre dans la zone de capture vidéo pour souffler des réponses.

Le second vecteur de la télésurveillance s'affranchit de l'analyse visuelle pour se concentrer sur l'empreinte comportementale numérique de l'apprenant, traquée silencieusement par les sondes intégrées dans l'interface de l'application cliente. Tout au long de l'épreuve chronométrée, des écouteurs d'événements enregistrent méticuleusement les tentatives de changement de fenêtre du navigateur, les excursions de la souris hors du périmètre de l'évaluation, ainsi que l'utilisation prohibée du presse-papiers pour coller des blocs de texte pré-rédigés. Toutefois, la véritable ingéniosité du système Fraudly ne réside pas dans la simple collection de ces événements épars, mais dans leur agrégation sémantique par un agent d'intégrité dédié. Cet algorithme de classification heuristique analyse les séquences temporelles des infractions détectées pour calculer un score de confiance dynamique. Une brève perte de focalisation oculaire due à un bruit extérieur ne déclenchera qu'une anomalie mineure, tandis qu'une série d'événements corrélés, associant la détection visuelle d'un smartphone suivie immédiatement d'une opération de collage de texte massif dans une question ouverte, sera interprétée comme un comportement frauduleux manifeste et caractérisé. À l'issue de l'examen, cet agent génère un rapport d'intégrité synthétique et probant, présentant à l'équipe pédagogique non seulement le verdict de suspicion de fraude, mais surtout la chronologie exacte des événements anormaux étayée par les captures d'écran décisives, fournissant ainsi les éléments factuels indispensables à toute prise de décision disciplinaire objective.
